%! Vectors — Unit 4, Lesson 4.8 — Homework
\documentclass[10pt]{article}
\usepackage{precalculus-article}
\usepackage{precalculus-boxes}
\newcommand{\TallMath}[1]{$\displaystyle #1\rule[-1.4em]{0pt}{3.2em}$}

\begin{document}

\pageheader{Unit 4, Lesson 4.8}{Homework}
\namedateperiod

\vspace{0.10in}

\begin{notesbox}{Vector Components and Magnitude}

\textbf{1.}\quad Find the component form and magnitude of the vector from
$P_1 = (-2, 1)$ to $P_2 = (4, 9)$.

\vspace{4pt}
$\vec{v} = \langle$ \blank{2cm}$,$ \blank{2cm}$\rangle$ \qquad
$\|\vec{v}\| = $ \blank{5cm}

\vspace{14pt}
\textbf{2.}\quad A force vector has magnitude $r = 13$ and direction angle
$\theta = 67.4^\circ$. Given $\cos 67.4^\circ \approx 0.385$ and
$\sin 67.4^\circ \approx 0.923$, find the component form of the vector.
Round to one decimal place.

\vspace{4pt}
$a = r\cos\theta \approx $ \blank{3.5cm} \qquad
$b = r\sin\theta \approx $ \blank{3.5cm}

\vspace{4pt}
$\vec{v} \approx \langle$ \blank{2cm}$,$ \blank{2cm}$\rangle$

\end{notesbox}

\vspace{0.08in}

\begin{notesbox}{Scalar Multiplication and Vector Addition}

\textbf{3.}\quad Let $\vec{u} = \langle 2, -3\rangle$ and $\vec{w} = \langle -4, 1\rangle$.
Compute each of the following.

\vspace{4pt}
(a)\quad $\vec{u} + \vec{w} = $ \blank{5cm}

\vspace{8pt}
(b)\quad $3\vec{u} = $ \blank{5cm}

\vspace{8pt}
(c)\quad $2\vec{u} - \vec{w} = \langle 2(2)-(-4),\; 2(-3)-1 \rangle = $ \blank{5cm}

\vspace{8pt}
(d)\quad Express $\vec{u}$ in terms of $\hat{\imath}$ and $\hat{\jmath}$:
$\vec{u} = $ \blank{5cm}

\end{notesbox}

\vspace{0.08in}

\begin{notesbox}{Unit Vectors and Dot Product}

\textbf{4.}\quad Find the unit vector in the direction of $\vec{v} = \langle -6, 8\rangle$.

\vspace{4pt}
$\|\vec{v}\| = $ \blank{3cm} \qquad
$\hat{u} = \left\langle \dfrac{-6}{\blank{1.2cm}},\; \dfrac{8}{\blank{1.2cm}}\right\rangle
= \langle$ \blank{2.5cm}$,$ \blank{2.5cm}$\rangle$

\vspace{14pt}
\textbf{5.}\quad Let $\vec{p} = \langle 3, -4\rangle$ and $\vec{q} = \langle 4, 3\rangle$.

\vspace{4pt}
(a)\quad Compute $\vec{p} \cdot \vec{q}$: \quad
$(3)(4) + (-4)(3) = $ \blank{4cm}

\vspace{8pt}
(b)\quad Are $\vec{p}$ and $\vec{q}$ perpendicular?
Circle one: \quad \textbf{Yes} \quad \textbf{No}

\vspace{4pt}
Justify: \writeline

\end{notesbox}

\vspace{0.08in}

\begin{notesbox}{Angle Between Vectors and Vector Addition Triangle}

\textbf{6.}\quad Two vectors have magnitudes $\|\vec{u}\| = 5$ and $\|\vec{w}\| = 7$,
and the angle between them is $\theta = 60^\circ$.

\vspace{4pt}
(a)\quad Use the Law of Cosines to find the magnitude of $\vec{r} = \vec{u} + \vec{w}$.
(Recall: $\|\vec{r}\|^2 = \|\vec{u}\|^2 + \|\vec{w}\|^2 - 2\|\vec{u}\|\,\|\vec{w}\|\cos\theta$
where $\theta$ is the angle between $\vec{u}$ and $\vec{w}$.)

\writelines{3}

\vspace{4pt}
(b)\quad Given $\vec{a} = \langle 1, 2\rangle$ and $\vec{b} = \langle 2, -1\rangle$,
use the dot product formula $\theta = \cos^{-1}\!\left(\dfrac{\vec{a}\cdot\vec{b}}{\|\vec{a}\|\,\|\vec{b}\|}\right)$
to find the angle between $\vec{a}$ and $\vec{b}$. Show all steps.

\writelines{4}

\end{notesbox}

\end{document}
